%% toc: P15 164 164.11
\section{Normal forms in the two model languages: proofs}
\label{sec:annex-p15-w02c-p15-028-model-languages}

The languages \(L_1,L_2\) and the counts \(m,m_U,l_U\) are those of
\cref{def:p15-w02c-p15-028-languages} and
\cref{def:p15-w02c-p15-028-counts}; \(D_v\), \(M_v\), \([U]\) and
\(\mathcal C(U)\) are those of \cref{def:p15-power-language}. Two facts
from \cref{thm:p15-power-normal-forms} and its proof are used
throughout: \(D_v\) is true at exactly the assignment \(v\), and
\(M_v\) is false at exactly the assignment \(v\); and for singular
sentences, agreement at every ground assignment is first-order
equivalence.

\begin{annexitem}{annex:p15-w02c-p15-028-l1-forms}{thm:p15-w02c-p15-028-l1-forms}{The \(L_1\) table}
\begin{proof}
\paragraph{The constituents.}
With \(A_1=P(a)\) and \(A_2=Q(a)\), the definition of \(M_v\) puts the
unnegated atom \(A_i\) into the disjunction when \(v_i=0\) and its
negation when \(v_i=1\). The four assignments \(00,01,10,11\) therefore
give the four displayed minimal sentences. The definition of \(D_v\)
puts \(A_i\) into the conjunction when \(v_i=1\) and \(\neg A_i\) when
\(v_i=0\), which gives the four displayed state descriptions. Since
\(n=2\), \(m=2^2=4\).

\paragraph{Truth sets.}
Evaluating each sentence at the four assignments gives
\[
\begin{aligned}
 [P(a)]&=\{11,10\}, & [Q(a)]&=\{11,01\},\\
 [P(a)\to Q(a)]&=\{11,01,00\}, & [P(a)\land Q(a)]&=\{11\},
\end{aligned}
\]
because \(P(a)\) is true when \(v_1=1\), \(Q(a)\) when \(v_2=1\), the
material conditional is false only at \(v_1=1,v_2=0\), and the
conjunction only at \(11\). A valid singular sentence is true at every
assignment, so its truth set is all of \(\Omega=\{00,01,10,11\}\). A
contradictory one is true at none, so its truth set is empty.

\paragraph{Normal forms.}
By \cref{thm:p15-power-normal-forms}, the maximal normal form of \(U\)
is the disjunction of the \(D_v\) with \(v\in[U]\), and the minimal
normal form is the conjunction of the \(M_v\) with \(v\notin[U]\); both
sets are unique. Reading off the truth sets:
\(P(a)\) has disjuncts \(D_{11},D_{10}\) and conjuncts
\(M_{00},M_{01}\); \(Q(a)\) has disjuncts \(D_{11},D_{01}\) and
conjuncts \(M_{00},M_{10}\); \(P(a)\to Q(a)\) has disjuncts
\(D_{11},D_{01},D_{00}\) and the single conjunct \(M_{10}\); a valid
sentence has all four disjuncts and no conjunct; a contradictory
sentence has no disjunct and all four conjuncts; \(P(a)\land Q(a)\) has
the single disjunct \(D_{11}\) and conjuncts \(M_{00},M_{01},M_{10}\).
The empty conjunction is \(\top\), equivalent to every valid sentence,
and the empty disjunction is \(\bot\), equivalent to every
contradictory sentence, as \cref{def:p15-power-language} stipulates.

The two forms the source writes out can also be checked directly by
the distributive and complement laws:
\[
\begin{aligned}
 (P(a)\lor Q(a))\land(P(a)\lor\neg Q(a))
   &\equiv P(a)\lor(Q(a)\land\neg Q(a))\equiv P(a),\\
 (P(a)\land Q(a))\lor(P(a)\land\neg Q(a))
   &\equiv P(a)\land(Q(a)\lor\neg Q(a))\equiv P(a),\\
 (P(a)\lor Q(a))\land(\neg P(a)\lor Q(a))
   &\equiv (P(a)\land\neg P(a))\lor Q(a)\equiv Q(a),
\end{aligned}
\]
and \(\neg P(a)\lor Q(a)\) is the material conditional
\(P(a)\to Q(a)\). For the three-disjunct form, the disjunction
\((P(a)\land Q(a))\lor(\neg P(a)\land Q(a))\) is equivalent to
\(Q(a)\), and \(Q(a)\lor(\neg P(a)\land\neg Q(a))\) is equivalent to
\((Q(a)\lor\neg P(a))\land(Q(a)\lor\neg Q(a))\), hence to
\(\neg P(a)\lor Q(a)\).

\paragraph{Counts.}
The columns \(m_U\) and \(l_U\) count the conjuncts and the
disjuncts listed above. In each row \(l_U+m_U=4=m\); this is the counting
identity \(|[U]|+|\mathcal C(U)|=2^n\) of
\cref{thm:p15-power-normal-forms}. For \(P(a)\), \(P(a)\to Q(a)\), a
valid and a contradictory sentence the values of \(m_U\) are
\(2,1,0,4\), and dividing by \(m=4\) gives \(g_1(U)=1/2,1/4,0,1\).

\paragraph{Union of content families.}
For singular \(S,S'\), the conjunction \(S\land S'\) is true at \(v\)
exactly when both conjuncts are, so \([S\land S']=[S]\cap[S']\), and
\[
 \Omega\setminus[S\land S']
   =(\Omega\setminus[S])\cup(\Omega\setminus[S']).
\]
Since \(\mathcal C(U)=\{M_v:v\in\Omega\setminus[U]\}\) and distinct
assignments index distinct minimal sentences,
\(\mathcal C(S\land S')=\mathcal C(S)\cup\mathcal C(S')\). With
\(S=P(a)\) and \(S'=Q(a)\) this is
\(\{M_{00},M_{01}\}\cup\{M_{00},M_{10}\}=\{M_{00},M_{01},M_{10}\}\),
which has \(2+2-1=3\) members because \(M_{00}\) is common to both. These
are the three minimal sentences \(P(a)\lor Q(a)\), \(P(a)\lor\neg Q(a)\)
and \(\neg P(a)\lor Q(a)\) that form \(K'\) in
\cref{thm:p15-power-source-examples}.
\end{proof}
\end{annexitem}

\begin{annexitem}{annex:p15-w02c-p15-028-l2-forms}{thm:p15-w02c-p15-028-l2-forms}{Counts and forms in \(L_2\)}
\begin{proof}
\paragraph{Atoms.}
A ground atom of \(L_2\) is a predicate applied to constants of the
right number. With the constants \(a,b\), the unary predicates \(P,Q\)
give \(2\cdot 2=4\) atoms, and the binary predicate \(R\) gives one
atom for each of the \(2^2=4\) ordered pairs of constants. There are no
function symbols and no identity, so there are no other atoms, and the
eight listed atoms are distinct expressions. Thus \(n=8\), and there are
\(2^8=256\) assignments. Distinct assignments give distinct state
descriptions and distinct minimal sentences, because \(D_v\) is true
and \(M_v\) false at \(v\) alone. Hence there are \(256\) of each.

\paragraph{The sample.}
Let \(v=00100011\). The definition of \(M_v\) puts in unnegated form
the atoms with \(v_i=0\) and negates the atoms with \(v_i=1\). These
are
\[
\begin{aligned}
 v_i=0:&\quad A_1,A_2,A_4,A_5,A_6
   =P(a),P(b),Q(b),R(a,a),R(a,b),\\
 v_i=1:&\quad A_3,A_7,A_8=Q(a),R(b,a),R(b,b).
\end{aligned}
\]
In the atom order this is literally the displayed disjunction. It is false
exactly when every disjunct is false, which is the stated assignment.

\paragraph{The forms of \(P(a)\) and \(Q(a)\).}
\(P(a)=A_1\) is false at \(v\) exactly when \(v_1=0\), so
\(\mathcal C(P(a))=\{M_v:v_1=0\}\). The remaining seven coordinates are
free, so this family has \(2^7=128\) members, and by
\cref{thm:p15-power-normal-forms} their conjunction is the unique
minimal normal form of \(P(a)\). A minimal sentence \(M_v\) contains
\(A_1\) unnegated exactly when \(v_1=0\), which identifies the
conjuncts as the minimal sentences with the unnegated disjunct
\(P(a)\). The same argument with the coordinate \(v_3\) of
\(Q(a)=A_3\) gives the \(128\)-member form of \(Q(a)\). The counting
identity gives \(l_U=256-128=128\), and \(g_1=128/256=1/2\).

\paragraph{Lifting the \(L_1\) forms.}
An \(L_1\) minimal sentence \(M^{L_1}_{v_1v_2}\) mentions only the
\(L_2\) atoms \(A_1=P(a)\) and \(A_3=Q(a)\), and it is false exactly
when \(P(a)\) has value \(v_1\) and \(Q(a)\) has value \(v_2\). In
\(L_2\) its truth set is therefore the complement of
\(\{w:w_1=v_1,\ w_3=v_2\}\), a set of \(2^6=64\) assignments, since the
other six coordinates are free. Its \(L_2\) content family is the
corresponding \(64\) minimal sentences \(M_w\). As \(64\neq 1\), it is
not itself an \(L_2\) minimal sentence.

The \(L_1\) minimal normal form of \(P(a)\) has the conjuncts
\(M^{L_1}_{00}\) and \(M^{L_1}_{01}\). Their \(L_2\) content families
are indexed by \(\{w:w_1=0,w_3=0\}\) and \(\{w:w_1=0,w_3=1\}\). These
sets are disjoint, since they differ at \(w_3\), and their union is
\(\{w:w_1=0\}\). By the union rule for content families proved above for
\cref{thm:p15-w02c-p15-028-l1-forms}, whose argument uses no feature
of \(L_1\) and so holds in any finite ground language, the union of the two content families is
\(\mathcal C(P(a))\) in \(L_2\), with \(64+64=128\) members. For
\(Q(a)\) the conjuncts are \(M^{L_1}_{00}\) and \(M^{L_1}_{10}\), with
index sets \(\{w:w_1=0,w_3=0\}\) and \(\{w:w_1=1,w_3=0\}\). These are
disjoint, and their union is \(\{w:w_3=0\}\).

\paragraph{Adding ground atoms in general.}
Let \(U\) be a singular sentence over \(A_1,\ldots,A_k\) and let the
larger language have atoms \(A_1,\ldots,A_n\). The truth value of \(U\)
at an assignment \(w\in\{0,1\}^n\) depends only on the restriction
\(w|_k=(w_1,\ldots,w_k)\). Hence
\[
 [U]_n=\{w:w|_k\in[U]_k\},
\]
and each \(u\in\{0,1\}^k\) has exactly \(2^{n-k}\) extensions. So
\(|[U]_n|=2^{n-k}|[U]_k|\) and, by the counting identity,
\(|\mathcal C_n(U)|=2^{n-k}|\mathcal C_k(U)|\). In the source's counts,
\(l_U\) and \(m_U\) are multiplied by \(2^{n-k}\), as is \(m=2^n\).
The ratio \(g_1(U)=m_U/m\) is unchanged. For \(P(a)\), passing from
\(L_1\) (\(k=2\)) to \(L_2\) (\(n=8\)), the factor is \(2^6=64\), and
\(m_{P(a)}\) goes from \(2\) to \(128\).
\end{proof}
\end{annexitem}

\begin{annexitem}{annex:p15-w02c-p15-028-nonsingular}{thm:p15-w02c-p15-028-nonsingular}{An essentially generalized consequence}
\begin{proof}
\paragraph{Item 1.}
Assume \(P(a)\lor Q(a)\) and, for reductio,
\(\forall x\,\neg(P(x)\lor Q(x))\). Universal instantiation at \(a\)
gives \(\neg(P(a)\lor Q(a))\), which contradicts the premise. Hence
\(\neg\forall x\,\neg(P(x)\lor Q(x))\), which is \(E\) with
\(\exists x\) read as \(\neg\forall x\neg\). Equivalently,
\(\exists\)-introduction from the instance at \(a\) gives \(E\) in one
step. Semantically, in any structure making \(P(a)\lor Q(a)\) true, the
denotation of \(a\) is a witness for \(E\).

\paragraph{Item 2.}
Let \(\mathfrak A\) have domain \(\{e\}\), interpret \(a\) as \(e\), and
let \(P\) and \(Q\) be empty. No element satisfies \(P(x)\lor Q(x)\),
so \(\mathfrak A\not\models E\). Thus \(E\) is not valid. By
\cref{thm:fol-soundness} it is not provable, so it is not
formally true in the sense of (7.1a).

\paragraph{Item 3.}
Let \(\mathfrak B\) have domain \(\{e,f\}\), interpret \(a\) as \(e\),
and put \(P=\{f\}\) and \(Q=\varnothing\). Then \(f\) witnesses
\(E\), so \(\mathfrak B\models E\). Both \(\mathfrak A\) and
\(\mathfrak B\) make \(P(a)\) and \(Q(a)\) false, so they have the same
ground assignment \(00\). A singular sentence of \(L_1\) is a Boolean
combination of \(P(a)\) and \(Q(a)\), so it has the same truth value in
\(\mathfrak A\) and \(\mathfrak B\). The sentence \(E\) does not. So no
singular sentence has the same truth value as \(E\) in every
structure. By \cref{thm:fol-soundness}, mutually derivable
sentences agree in every structure, so \(E\) is not equivalent in the
sense of (7.1a) to any singular sentence. This is what the source calls
essentially generalized.

\paragraph{Item 4.}
This is item 3 for the singular sentence \(P(a)\lor Q(a)\). Directly,
\(\mathfrak B\) satisfies \(E\) and falsifies \(P(a)\lor Q(a)\). So
\(E\) does not entail \(P(a)\lor Q(a)\), and the two are not
equivalent. Items 1, 2 and 4 together show that \(E\) is a consequence
of the minimal sentence \(M_{00}\) that is neither equivalent to
\(M_{00}\) nor valid. Hence the singular hypothesis in the clause of
\cref{thm:p15-power-normal-forms} cannot be removed.

\paragraph{Mutual exclusiveness for arbitrary consequences.}
Let \(v\neq w\), and let \(H\) be any sentence with
\(M_v\entails H\) and \(M_w\entails H\). By
\cref{thm:p15-power-normal-forms}, \(M_v\lor M_w\) is valid. So every
structure satisfies \(M_v\) or \(M_w\), and in either case it satisfies
\(H\). Hence \(H\) is valid. No singularity of \(H\) was used.

\paragraph{The dual boundary.}
Let \(F=D_{11}\land\forall x\,P(x)\). Then \(F\entails D_{11}\) by
conjunction elimination. The structure with domain \(\{e\}\), \(a\)
interpreted as \(e\) and \(P=Q=\{e\}\) satisfies \(F\), so \(F\) is
consistent. The structure with domain \(\{e,f\}\), \(a\) interpreted
as \(e\), \(P=\{e\}\) and \(Q=\{e\}\) satisfies \(D_{11}\) but not
\(\forall x\,P(x)\), hence not \(F\). Thus \(D_{11}\not\entails F\),
and \(F\) is a consistent, non-singular sentence that entails
\(D_{11}\) without being equivalent to it.
\end{proof}
\end{annexitem}
